\section{Con số, và điều nó không đếm}
\label{sec:ch12-con-so-va-cai-no-dem}

\index{kiểm đếm tài liệu XAI!bốn mức lọc}%
Với thủ tục tìm và thủ tục chấm của mục trước, kết quả là một cái phễu bốn
mức, và bảng 3 của bài báo in cả bốn~\autocite{p24humangap}. Mức thứ nhất là
toàn bộ tài liệu mà truy vấn gom được: 18\,254 bài mang từ khóa về khả năng
giải thích. Mức thứ hai lọc lấy những bài chứa cả từ khóa về thử nghiệm với
người: 253 bài, tức 1.39\%. Mức thứ ba loại 16 bài được thống nhất là lạc
đề, tức những bài không trình bày một phương pháp XAI nào mà là bài rà tổng
quan về chính lĩnh vực XAI, bài bàn triết học về khả năng giải thích, hoặc bài
tóm lược công trình của người khác; còn lại 237 bài đúng chủ đề, tức 1.30\%. Mức thứ tư là những
bài thật sự có kiểm chứng với người: 128 bài, tức 0.70\%.

Bốn con số ấy khớp với nhau ở hai chỗ mà một cái phễu có thể lệch. Bảng 2 của
bài chấm 253 bài của mức hai thành 109 bài không có kiểm chứng, 128 bài có, và
16 bài lạc đề; ba số ấy cộng lại đúng 253, và 253 trừ 16 đúng bằng 237 của mức
ba~\autocite{p24humangap}.

\index{kiểm đếm tài liệu XAI!tỉ lệ 0.7 phần trăm}%
Con số 0.70\% là con số của nhan đề, và nó là 128 trên 18\,254. Đọc nó theo
chiều ngược lại thì bài báo nói rằng hơn 98\% tài liệu XAI thậm chí không xuất
hiện cùng từ khóa nào về đánh giá với người, và hơn 99\% không chứa bằng chứng
thực nghiệm nào từ người~\autocite{p24humangap}. Các tác giả cũng nói con số
này không nhạy với sai số của bộ từ khóa theo cách người ta có thể nghĩ: giả
sử tiêu chí tìm sai tới một bậc độ lớn, tức tỉ lệ kiểm chứng thật cao gấp mười
lần, thì kết quả vẫn chỉ là 7\% số công trình được thử với
người~\autocite{p24humangap}.

Trong bốn con số, con số đáng dừng lại lâu nhất là tỉ lệ ở mức ba so với mức
bốn, tức 128 trong 237. Hai mức ấy nói về cùng một tập bài: những bài đã tự
tuyên bố rằng phương pháp của mình giải thích được cho người. Vậy trong số
chính những bài đã hứa điều đó, chỉ hơn một nửa đem lời hứa ra thử, và gần một
nửa thì không.
Bài báo phát biểu tỉ lệ này hai lần và hai lần không khớp nhau: mục 3 viết
\enquote{54\% of papers provide an evaluation, 46\% do not}, còn mục 4.1 viết
\enquote{only 56\% of them provided an evaluation with
humans}~\autocite{p24humangap}. Phép chia 128 cho 237 cho 54\%, và cả hai bảng
đều nhất quán với mục 3, nên con số ở mục 4.1 là chỗ sai của bài. Sách vì thế
in các phép đếm chứ không in tỉ lệ ấy, và ghi lại chỗ vênh này một lần tại đây;
nhận xét này là của sách.

\index{kiểm đếm tài liệu XAI!ba khảo sát trước đó}%
Ba công trình trước đó mà bài báo dựa lên cùng cho một hình dạng, và sách thuật
lại cả ba qua lời của bài báo vì không đọc bài nào trong ba. Miller và cộng sự
năm 2017 rà 23 bài của một hội thảo XAI: chỉ 4 bài dẫn tới nghiên cứu khoa học
xã hội liên quan, chỉ 1 bài dựng mô hình trên nền ấy, và không bài nào trong
23 có kiểm chứng thực nghiệm. Wells và Bednarz rà 25 bài về giải thích cho học
tăng cường và thấy 17 bài không có thử nghiệm nào với người dùng. Nauta và
cộng sự rà 300 bài ở các hội nghị AI hàng đầu và thấy cứ 3 bài thì 1 bài đánh
giá phương pháp của mình chỉ bằng bằng chứng giai thoại, còn cứ 5 bài mới có 1
bài đánh giá với người dùng thật~\autocite{p24humangap}. Ba cuộc rà ấy đọc 23,
25 và 300 bài, so với 18\,254 bài của cuộc kiểm đếm này, và ba lần đều ra cùng
một kết luận.

\index{độ trung thực!vắng tên trong cuộc kiểm đếm}%
Đến đây phải nói rõ điều mà con số 0.70\% không nói, vì đây là chỗ dễ đọc lệch
nhất trong cả chương. Trong toàn bộ bài báo, các chữ \enquote{faithfulness} và
\enquote{faithful} không xuất hiện lần nào, chữ \enquote{ground truth} cũng
không, và cả \enquote{attribution}, \enquote{LIME}, \enquote{SHAP} lẫn
\enquote{saliency} đều không~\autocite{p24humangap}. Cái được đếm là việc kiểm
chứng lời tuyên bố về khả năng giải thích cho người. Theo cặp khái niệm mà
mục~\ref{sec:ch01-do-trung-thuc-va-tinh-hop-ly} tách ra, đó là phía tính hợp
lý, không phải phía độ trung thực. Vậy nên câu \enquote{chưa tới 1\% số bài XAI
kiểm chứng độ trung thực} là một câu sai, và không phải câu bài báo này nói.

Nếu vậy thì con số ấy nói gì với Phần~IV? Nó nói đúng điều mà quan hệ vay mượn
ở mục~\ref{sec:ch12-hai-nhanh-danh-gia} đòi. Đánh giá theo chức năng đứng vững
nhờ một việc kiểm chứng đã làm ở cách đánh giá đắt hơn, và cuộc kiểm đếm cho
thấy việc kiểm chứng ấy hầu như không được làm ở đâu cả. Món nợ không phải là
món nợ đã trả rồi mà ta quên mất; nó chưa từng được trả. Đó là lý do một chương
về vắng mặt con người thuộc về Phần~IV chứ không thuộc về Phần~V: nó không đề
xuất lối thoát nào, nó ghi lại rằng cái nền mà cả Phần~IV đứng lên chưa được ai
dựng.

Thêm một điều nữa mà phép đếm không nói, và bài báo tự nêu nó. Một bài được 1
điểm khi nó có một thí nghiệm hành vi với người, bất kể thí nghiệm ấy tốt hay
tồi. Các tác giả viết rằng thang điểm này có thể cần làm lại, vì sự có mặt đơn
thuần của một cuộc đánh giá chưa chắc tương đương với
\enquote{a paper that presents a \enquote{good} evaluation}, và họ dẫn ngay quan sát
của chính mình rằng một số đánh giá không ghi cả số người tham gia lẫn chuyên
môn của họ~\autocite{p24humangap}. Nói cách khác, 128 là chặn trên chứ không phải chặn
dưới. Mục~\ref{sec:ch12-nghien-cuu-khong-co-nguoi} đọc chỗ đó.

\index{kiểm đếm tài liệu XAI|)}%
